\chapter{ToyMCMC: Independent-Chain Random-Walk MCMC}
\label{ch:toymcmc}

\newthought{This chapter is for anyone who wants a posterior-shaped point cloud} ---
not the Bayesian evidence Chapters~\ref{ch:multinest}--\ref{ch:dynesty} compute, just
the shape of the posterior itself, at a fraction of nested sampling's per-point
bookkeeping. It is also the chapter that teaches the \textsc{mcmc} family's runtime:
how \Jtwo\ drives many Markov chains through the distributed Worker fleet, using
\sampler{ToyMCMC} --- the plain random-walk member of the family --- as the vehicle.
Chapter~\ref{ch:ptmcmc} builds directly on everything here and covers only what
parallel tempering adds. After this chapter you can explain what a Metropolis chain is
actually doing, read its \textsc{yaml} block fluently, and recognise every line it
prints while running.

\section{What a Metropolis chain computes}
\label{sec:mcmc-what}

Nested sampling (Chapters~\ref{ch:multinest}--\ref{ch:dynesty}) shrinks a population
through the likelihood to compute an integral. \textsc{mcmc} does something simpler:
it constructs a random walk whose \emph{stationary distribution} --- the distribution
its visited points converge to, given enough steps --- is exactly the (unnormalised)
posterior. It never estimates the normalising evidence; it only needs the
likelihood \emph{ratio} between two points, so the constant you would need for
$\log Z$ cancels out and is never computed.

\subsection{The algorithm: Metropolis--Hastings}
\label{sec:mcmc-algorithm}

Each chain holds one current position $u$ in the unit cube. One step is:

\begin{enumerate}
  \item \textbf{Propose.} Draw $u' = u + \delta$, where $\delta \sim
        \mathcal{N}(0, \yamlkey{proposal\_scale}^2)$ componentwise --- a symmetric
        random walk\cite{Metropolis:1953,Hastings:1970}. If $u'$ would land outside
        $[0,1]^d$, redraw $\delta$ and try again; the step is never clipped to the
        boundary, only redrawn until it lands inside.
  \item \textbf{Evaluate.} Send $u'$ through the full declared pipeline
        (Chapters~\ref{ch:calculators}--\ref{ch:loglikelihood}) to get its
        log-likelihood $\log L'$.
  \item \textbf{Accept or reject.} Let $\Delta = \beta\,(\log L' - \log L)$, with
        $\beta = 1$ for \sampler{ToyMCMC} (temperature is always $1$; $\beta$ becomes
        the inverse temperature once Chapter~\ref{ch:ptmcmc} introduces a ladder).
        Accept ($u \leftarrow u'$) if $\Delta \geq 0$, or with probability
        $e^{\Delta}$ otherwise; keep $u$ unchanged if rejected.
  \item \textbf{Repeat} for \yamlkey{num\_iters} steps.
\end{enumerate}

A chain's very first proposal is always accepted (there is no previous likelihood to
compare against yet); every later step follows the rule above. Why a rule this simple
is enough to guarantee the chain's stationary distribution is the posterior is worth
spelling out once.

\subsection{Why the rule works: detailed balance}
\label{sec:mcmc-detailed-balance}

A Markov chain is fully described by a transition kernel $T(u \to u')$: the
probability of moving to $u'$ given the chain currently sits at $u$. Under mild
conditions --- the chain can reach anywhere from anywhere, and never cycles
deterministically --- repeatedly applying $T$ from any starting point converges to a
unique \emph{stationary distribution}: a distribution $\pi$ that, once reached, $T$
leaves unchanged. This is the formal justification for treating a long enough
chain's visited points as draws from $\pi$ --- it is the distribution a chain forgets
its own starting point in favour of.

Constructing a $T$ whose stationary distribution is a specific $\pi$ chosen in
advance is generally hard to do directly. The standard shortcut is \emph{detailed
balance} (or reversibility): a kernel satisfying $\pi(u)\,T(u \to u') = \pi(u')\,T(u'
\to u)$ for every pair $u, u'$ is guaranteed to have $\pi$ as a stationary
distribution. It is a sufficient, not necessary, condition, but the one nearly every
practical MCMC method is actually built around, because checking it only ever
requires reasoning about one pair of states at a time\cite{RobertCasella:2004}.

Metropolis--Hastings builds a detailed-balanced $T$ from two pieces: a proposal
density $q(u' \mid u)$, chosen freely (here, the Gaussian random walk in step~1
above), and an acceptance probability $\alpha(u \to u') = \min\big(1,\ \pi(u')\,q(u
\mid u')\,/\,[\pi(u)\,q(u' \mid u)]\big)$, derived by requiring the resulting kernel
to satisfy detailed balance exactly\cite{Metropolis:1953,Hastings:1970}. Only the
ratio $\pi(u')/\pi(u)$ survives inside $\alpha$ --- which is why an
\emph{unnormalised} target is enough. The posterior here is $\pi(u) \propto L(u)$
under the flat prior every \yamlkey{Bounds} block defines on $[0,1]^d$, so a
Metropolis chain only ever compares two likelihoods and never needs the normalising
evidence integral that Chapters~\ref{ch:multinest}--\ref{ch:dynesty} exist to
compute.

For a \emph{symmetric} proposal --- $q(u' \mid u) = q(u \mid u')$, true of the
Gaussian random walk here, since moving from $u$ to $u'$ is exactly as likely, before
either is evaluated, as the reverse move --- the proposal ratio is $1$ and $\alpha$
collapses to $\min(1,\ \pi(u')/\pi(u))$: exactly step~3 of the algorithm above, with
$q$ never appearing at all. This symmetric case is Metropolis, Rosenbluth,
Rosenbluth, Teller \& Teller's original algorithm\cite{Metropolis:1953}; Hastings's
contribution\cite{Hastings:1970} was generalising it to \emph{asymmetric} proposals.
\sampler{ToyMCMC}'s proposal is symmetric by construction, which is the actual reason
\S\ref{sec:mcmc-algorithm} never needed to mention $q$ at all.

\begin{jtip}
A random walk with too small a \yamlkey{proposal\_scale} accepts almost everything
but barely moves; too large, and it rejects almost everything. For a random-walk
Metropolis chain in high dimensions, the acceptance rate that maximises how fast the
chain explores per unit of computation converges to $\approx 23.4\,\%$ as the number
of parameters grows\cite{RobertsGelmanGilks:1997} --- in practice, most
practitioners tune toward an acceptance rate somewhere around 20--40\,\%
(\S\ref{sec:toymcmc-yaml}), neither extreme, and a useful first diagnostic when a run
looks unhealthy.
\end{jtip}

\section{How the code implements it}
\label{sec:toymcmc-code}

\subsection{Class hierarchy: one engine, thin method subclasses}
\label{sec:toymcmc-classes}

\begin{quote}
\sffamily\footnotesize
\code{FeedbackSampler} $\;\to\;$ \code{MCMCBaseSampler} $\;\to\;$ \code{ToyMCMCSampler}
\end{quote}

\code{FeedbackSampler} is the same shared base nested sampling and
\sampler{AdaptiveBridson} build on (Chapters~\ref{ch:multinest},
\ref{ch:adaptive-levelset}) --- checkpointing and the notion of a barrier.
\code{MCMCBaseSampler} adds everything generic to the whole \textsc{mcmc} family: the
chain registry, the Redis transport (\S\ref{sec:mcmc-runtime}), deterministic
\textsc{uuid}s, checkpoint export/import, and the accept/reject bookkeeping.
\code{ToyMCMCSampler} itself is a thin subclass: it sets \code{method = "ToyMCMC"},
normalises \yamlkey{proposal\_scale} into one value per chain, and attaches a progress
observer --- nothing about the Markov-chain science lives here, because plain
random-walk Metropolis \emph{is} what \code{MCMCBaseSampler} already implements by
default.

\begin{jnote}
This is why \sampler{ToyMCMC} is the vehicle for this chapter rather than an
afterthought: every other independent-chain method in the family
(Section~\ref{sec:mcmc-family} in Chapter~\ref{ch:choosing-sampler}) is the same base
class with a different chain engine substituted in --- learning \sampler{ToyMCMC}
teaches the runtime every one of them uses.
\end{jnote}

\subsection{The chain engine}
\label{sec:toymcmc-engine}

Each chain owns one \code{MCMCChain} instance --- current parameter vector $u$,
\yamlkey{proposal\_scale}, an iteration counter, and its own \code{numpy} random
generator. Chain RNGs are never shared: one child \code{SeedSequence} is spawned per
chain from the card's \yamlkey{Bounds.seed}, so a chain's trajectory depends only on
\code{(seed, chain\_id)} --- never on how many Workers happened to be running or in
what order their results arrived.

Every proposed point gets a deterministic \textsc{uuid} minted from
\code{(method, seed, chain\_id, step, stage)} through a \textsc{sha}-256 digest ---
the same contract Chapter~\ref{ch:choosing-sampler}'s reproducibility tip promised for
the whole fixed-set/adaptive/\textsc{mcmc} side of the book (unlike nested sampling,
whose proposal \textsc{uuid}s are fresh every time). Fix \yamlkey{seed} and a
\sampler{ToyMCMC} run is byte-for-byte re-proposable.

\subsection{The runtime: async independent chains}
\label{sec:mcmc-runtime}

Chains that share no cross-chain move --- \sampler{ToyMCMC} among them --- run on
\Jtwo's \emph{async independent-chain pipeline}, the base design every non-coupled
\textsc{mcmc} method uses:

\begin{enumerate}
  \item \textbf{One shared task queue.} Every chain's proposal lands on the same
        Redis list, \code{hep:task\_queue} --- Workers steal freely, exactly as for
        any other method (Section~\ref{sec:redis-schema}).
  \item \textbf{Per-chain feedback shards.} Results come back on
        \code{hep:feedback:chain:\{id\}}, one list per chain, not the single shared
        \code{hep:feedback} list other methods use --- so the control process can
        wait for \emph{any} chain's next result without being blocked behind an
        unrelated chain's slower one.
  \item \textbf{One in-flight proposal per chain.} At most one outstanding evaluation
        exists per chain at any time. The moment a chain's result comes back, it is
        absorbed and that chain is immediately re-proposed --- there is no
        cross-chain generation barrier, so a fast chain can run many steps ahead of a
        slow one.
\end{enumerate}

\begin{jnote}
This is the opposite scheduling shape from nested sampling
(\S\ref{sec:nested-distributed}): there, one single decision (which point to retire
next) gates the whole population. Here, $N$ independent decisions run fully in
parallel, each gated only by its own chain's last result --- so
\sampler{ToyMCMC} keeps scaling with more Workers for longer than a single-threaded
nested-sampling control loop does.
\end{jnote}

\begin{jwarn}
Chapter~\ref{ch:ptmcmc}'s \sampler{PTMCMC} does \emph{not} use this pipeline ---
temperature exchange needs every replica to reach the same point at the same time, so
coupled \textsc{mcmc} methods fall back to the classical generation-barrier path
instead (\S\ref{sec:ptmcmc-runtime}). Run \cli{Jarvis man sampler.mcmc-runtime} for
both pipelines side by side.
\end{jwarn}

Prefer \yamlkey{num\_chains} at least as large as \yamlkey{EnvReqs.V2.workers}
(Chapter~\ref{ch:envreqs}) --- with fewer chains than Workers, some Workers have
nothing to steal and idle. \Jtwo\ warns exactly this at startup when it detects the
mismatch.

\section{Configuring ToyMCMC}
\label{sec:toymcmc-yaml}

The minimal working card:

\begin{yamlwin}
Sampling:
  Method: ToyMCMC
  Variables:
    - name: x
      distribution: {type: Flat, parameters: {min: 0.0, max: 1.0}}
    - name: y
      distribution: {type: Flat, parameters: {min: 0.0, max: 1.0}}
  Bounds:
    num_chains: 4
    num_iters: 2000
    proposal_scale: 0.2
    seed: 0
  LogLikelihood:
    - {name: LogL, expression: "..."}   # your physics
\end{yamlwin}

Table~\ref{tab:toymcmc-bounds} lists every \yamlkey{Bounds} key.

\begin{table}[ht]
  \footnotesize
  \begingroup
  \setlength{\tabcolsep}{0pt}
  \begin{tabular}{@{}p{0.229\linewidth}p{0.187\linewidth}p{0.584\linewidth}@{}}
    \toprule
    \textbf{Key} & \textbf{Default} & \textbf{Meaning} \\
    \midrule
    \yamlkey{num\_chains} & \req & how many independent chains run (prefer
      $\geq$ \yamlkey{workers}, \S\ref{sec:mcmc-runtime}) \\
    \yamlkey{num\_iters} & \req & Metropolis transitions per chain; the run submits
      exactly \yamlkey{num\_chains}$\,\times\,$\yamlkey{num\_iters} proposals in total \\
    \yamlkey{proposal\_scale} & \req & the random-walk step's standard deviation
      (\S\ref{sec:mcmc-algorithm}) --- a scalar or a list \\
    \yamlkey{seed} & \code{0} & master \textsc{rng} seed; drives every chain's own
      child seed and every proposal's \textsc{uuid} \\
    \bottomrule
  \end{tabular}
  \endgroup
  \caption{\sampler{ToyMCMC} \yamlkey{Bounds} keys. \code{Jarvis man sampler.ToyMCMC}
  is the live reference.}
  \label{tab:toymcmc-bounds}
\end{table}

\begin{keylist}
  \item[proposal\_scale] One positive number broadcasts to every chain; a one-element
        list broadcasts the same way. A list of exactly \yamlkey{num\_chains} values
        gives each chain its own step size --- useful once you know some regions of
        the posterior need a smaller step than others, without running separate
        cards.
\end{keylist}

\begin{jnote}
Every sampler-specific setting lives under \yamlkey{Bounds}, never at the
\yamlkey{Sampling} top level --- the same contract nested sampling and every other
method share (Chapter~\ref{ch:choosing-sampler}).
\end{jnote}

\section{Validating your card}
\label{sec:toymcmc-validation}

\begin{terminal}
\begin{jcmd}
Jarvis validate scan.yaml
\end{jcmd}
\end{terminal}

An empty \yamlkey{Bounds} block is rejected before anything runs. Each missing key is
reported twice --- once by the schema, once by the method's own check --- so three
omissions make six findings, and any report of two or more opens with a summary table
before the detail:

\begin{terminal}
\begin{jout}
Config validation failed (6 error(s), 0 warning(s)):

Error summary (reference only):
+---+-------------+--------------------------------+------------------------+
| # | Code        | YAML path                      | Problem                |
+---+-------------+--------------------------------+------------------------+
| 1 | JV2-SCH-001 | \$.Sampling.Bounds              | 'num_chains' is a r... |
| 2 | JV2-SCH-001 | \$.Sampling.Bounds              | 'num_iters' is a re... |
| 3 | JV2-SCH-001 | \$.Sampling.Bounds              | 'proposal_scale' is... |
| 4 | JV2-MTH-051 | Sampling.Bounds.num_chains     | ToyMCMC requires Sa... |
| 5 | JV2-MTH-052 | Sampling.Bounds.num_iters      | ToyMCMC requires Sa... |
| 6 | JV2-MTH-053 | Sampling.Bounds.proposal_scale | ToyMCMC requires Sa... |
+---+-------------+--------------------------------+------------------------+
This table is only a quick reference. Please consult the Jarvis-HEP YAML settings
documentation before editing the task card.

Detailed diagnostics:
  [error] JV2-SCH-001  \$.Sampling.Bounds
          'num_chains' is a required property
          hint: Run: Jarvis man yaml.Sampling.Bounds
          suggestion: Add 'num_chains' at this location using the expected
          YAML type.
  ...
  [error] JV2-MTH-051  Sampling.Bounds.num_chains
          ToyMCMC requires Sampling.Bounds.num_chains (a positive integer)
          hint: Run: Jarvis man sampler.ToyMCMC
          suggestion: Set the ToyMCMC chain count and iteration count to
          positive integers.
          example:
            Bounds:
              num_chains: 4
              num_iters: 2000
\end{jout}
\end{terminal}

Table~\ref{tab:toymcmc-validation} lists the codes a \sampler{ToyMCMC} card is most
likely to trigger.

\begin{table}[ht]
  \footnotesize
  \begingroup
  \setlength{\tabcolsep}{0pt}
  \begin{tabular}{@{}p{0.171\linewidth}p{0.297\linewidth}p{0.532\linewidth}@{}}
    \toprule
    \textbf{Code} & \textbf{Fires on} & \textbf{Fix} \\
    \midrule
    \code{JV2-MTH-050} & \yamlkey{Bounds} missing entirely & add \yamlkey{num\_chains},
      \yamlkey{num\_iters}, \yamlkey{proposal\_scale} \\
    \code{JV2-MTH-051} & \yamlkey{num\_chains} missing or $< 1$ & set a positive
      integer \\
    \code{JV2-MTH-052} & \yamlkey{num\_iters} missing or $< 1$ & set a positive
      integer \\
    \code{JV2-MTH-053} & \yamlkey{proposal\_scale} missing, non-positive, or the
      wrong length & one positive number, or one per chain \\
    \code{JV2-MTH-055} & \yamlkey{proposal\_scale} list length is neither $1$ nor
      \yamlkey{num\_chains} & match the length or use a scalar \\
    \code{JV2-MTH-056} & \yamlkey{seed} present but negative or non-integer & use a
      non-negative integer \\
    \bottomrule
  \end{tabular}
  \endgroup
  \caption{Validation codes a \sampler{ToyMCMC} card is most likely to hit. The other
  methods' codes are in Appendix~\ref{app:troubleshooting}.}
  \label{tab:toymcmc-validation}
\end{table}

\section{Reading the log}
\label{sec:mcmc-log}

Every excerpt below is a real \cli{Jarvis run} transcript against a two-parameter
Gaussian bump, \yamlkey{num\_chains: 4}, \yamlkey{num\_iters: 50},
\yamlkey{proposal\_scale: 0.2}. Everything lands in \file{sampler.log}, using the
visual format of Section~\ref{sec:log-visual-format}.

\subsection{At startup}
\label{sec:toymcmc-log-startup}

\begin{terminal}[sampler.log]
\begin{jout}
·•· Jarvis-HEP.Sampler.Toymcmc
	-> 08-13 00:14:39.833 - [WARNING] >>>
Initializing the ToyMCMC Sampling

·•· Jarvis-HEP.Sampler.Toymcmc
	-> 08-13 00:14:39.834 - [INFO] >>>
ToyMCMC Sampler configured: chains=4 iterations=50 total_transitions=200
proposal_scale=0.2

·•· Jarvis-HEP.Sampler.Toymcmc
	-> 08-13 00:14:39.834 - [INFO] >>>
ToyMCMC independent-chain pipeline: chains=4 iters=50 sharded_feedback=True
\end{jout}
\end{terminal}

\code{total\_transitions} is \yamlkey{num\_chains}$\,\times\,$\yamlkey{num\_iters} ---
the exact number of Metropolis steps (not accepted moves) the run will submit before
it is done.

\subsection{While it runs: the progress meter}
\label{sec:toymcmc-log-progress}

Every absorbed transition advances a per-mille progress meter, following the same
convention as nested sampling's every-100th-iteration escalation
(Section~\ref{sec:nested-log-progress}): an \code{[INFO]} line every $1/1000$ of the
total, promoted to \code{[WARNING]} every whole percent so long runs stay visible in a
busy log without scrolling:

\begin{terminal}[sampler.log]
\begin{jout}
·•· Jarvis-HEP.Sampler.Toymcmc
	-> 08-13 00:14:40.279 - [WARNING] >>>
10‰ of 2/200 ToyMCMC transitions completed in 00:00:00.445
(chains=4 iterations=1/50)
\end{jout}
\end{terminal}

\code{2/200} is transitions absorbed so far out of the total; \code{iterations=1/50}
is the \emph{furthest-ahead} chain's own step count --- with async independent chains,
different chains are rarely at exactly the same iteration.

\subsection{At the end: the summary}
\label{sec:toymcmc-log-summary}

\begin{terminal}[sampler.log]
\begin{jout}
·•· Jarvis-HEP.Sampler.Toymcmc
	-> 08-13 00:14:40.655 - [INFO] >>>
ToyMCMC diagnostics sampler_summary -> DATABASE/sampler_summary.json

·•· Jarvis-HEP.Sampler.Toymcmc
	-> 08-13 00:14:40.655 - [INFO] >>>
ToyMCMC diagnostics chain_history -> DATABASE/chain_history.csv

·•· Jarvis-HEP.Sampler.Toymcmc
	-> 08-13 00:14:40.709 - [INFO] >>>
ToyMCMC finished: proposed=200 accepted=60 accept_rate=0.3000 async=True
\end{jout}
\end{terminal}

This run's $30\,\%$ acceptance rate is squarely in the healthy range
(\S\ref{sec:mcmc-algorithm}'s tip) for a $0.2$ proposal scale on a unit-cube Gaussian
bump.

\section{What the run leaves behind}
\label{sec:mcmc-outputs}

Alongside the normal output tree (Chapter~\ref{ch:outputs}), every \textsc{mcmc}-family
run writes two extra files under \file{DATABASE/}, in addition to the usual
\file{samples.csv} / \file{samples.hdf5}:

\begin{table}[ht]
  \footnotesize
  \begingroup
  \setlength{\tabcolsep}{0pt}
  \begin{tabular}{@{}p{0.255\linewidth}p{0.745\linewidth}@{}}
    \toprule
    \textbf{File} & \textbf{Contains} \\
    \midrule
    \code{sampler\_summary.json} & one record: \yamlkey{accept\_rate},
      \yamlkey{rhat\_logl}, \yamlkey{ess\_logl\_mean}, \yamlkey{failed\_samples}, and
      a per-chain breakdown (\S\ref{sec:mcmc-diagnostics}) \\
    \code{chain\_history.csv} & one row per absorbed transition: \yamlkey{chain\_id},
      \yamlkey{step}, \yamlkey{accepted}, \yamlkey{weight}, \yamlkey{logl},
      \yamlkey{temperature}, \yamlkey{state} \\
    \bottomrule
  \end{tabular}
  \endgroup
  \caption{\textsc{mcmc}-family diagnostics, additive to \code{samples.hdf5} /
  \code{samples.csv}.}
  \label{tab:mcmc-outputs}
\end{table}

\begin{jnote}
\textbf{Every proposal is archived, accepted or not.} Unlike nested sampling, where
\code{ncall} usually exceeds \code{niter} because rejected candidates are never kept
(Section~\ref{sec:nested-outputs}), every \textsc{mcmc} proposal is a fully evaluated,
fully archived Sample --- \code{total\_proposed} in \file{sampler\_summary.json}
\emph{is} the row count in \file{samples.csv} for this scan. Rejection only decides
which point the chain's \emph{next} proposal builds from; it never discards a
computed point from the archive.
\end{jnote}

\yamlkey{chain\_id}, \yamlkey{step}, and \yamlkey{stage} are also stamped directly into
every Sample's own observables, so they appear as ordinary columns in
\file{samples.csv} / \file{samples.hdf5} alongside your physical parameters --- group
by \yamlkey{chain\_id} to recover each chain's own trajectory from the plain archive,
with no need to touch \file{chain\_history.csv} unless you want the accept/reject
ledger specifically.

\section{Resume}
\label{sec:mcmc-resume}

Unlike nested sampling, which checkpoints only once at the very end
(\S\ref{sec:nested-distributed}), the \textsc{mcmc} family checkpoints
\emph{during} the run --- time-cadenced, by default every $30$ seconds, at whichever
feedback boundary the method's pipeline allows (a single absorbed transition for
async independent chains; a full generation for the barrier-coupled family in
Chapter~\ref{ch:ptmcmc}). An interrupted run leaves a real, resumable snapshot:

\begin{terminal}[real transcript]
\begin{jout}
·•· Jarvis-HEP
	-> 00:16:20.478 - [WARNING] >>>
received SIGINT — initiating clean shutdown (Workers / Archiver / managed Redis)

·•· Jarvis-HEP
	-> 00:16:20.481 - [WARNING] >>>
interrupt checkpoint written -> checkpoints/<scan>/ToyMCMC/state.pkl
(resume: Jarvis run <task.yaml> --resume)
\end{jout}
\end{terminal}

\begin{terminal}
\begin{jcmd}
Jarvis run scan.yaml --resume
\end{jcmd}
\end{terminal}

\begin{terminal}[real transcript]
\begin{jout}
·•· Jarvis-HEP
	-> 00:16:46.048 - [WARNING] >>>
Resume DATABASE reconciliation found prefix=0, completed=4 failed=0, 4 legacy UUID(s)

·•· Jarvis-HEP
	-> 00:16:46.049 - [WARNING] >>>
Resume checkpoint path -> checkpoints/<scan>/ToyMCMC/state.pkl (payload=yes)
\end{jout}
\end{terminal}

\begin{jwarn}
The checkpoint is validated against the \emph{current} card before it is trusted: if
\yamlkey{num\_chains}, \yamlkey{num\_iters}, \yamlkey{proposal\_scale}, or
\yamlkey{seed} changed since the checkpoint was written, resume refuses rather than
silently continuing under a shape the summary no longer describes --- edit any of
those and the run must restart clean.
\end{jwarn}

\section{Diagnostics: R-hat and effective sample size}
\label{sec:mcmc-diagnostics}

\file{sampler\_summary.json} carries two convergence diagnostics, computed on each
chain's accepted log-likelihood trace (discarding the first half as burn-in):

\begin{keylist}
  \item[rhat\_logl] The Gelman--Rubin $\hat{R}$ statistic\cite{GelmanRubin:1992}:
        compares between-chain and
        within-chain variance of the log-likelihood trace. Values near $1.0$ suggest
        the chains agree on where the posterior mass is; a real run of this chapter's
        example card gave $\hat{R}=1.044$ across four chains.
  \item[ess\_logl\_mean] Mean effective sample size across chains, from the
        log-likelihood trace's autocorrelation time --- how many effectively
        independent draws the accepted points represent, which is always smaller than
        the raw accepted count for a correlated random walk.
\end{keylist}

\begin{jwarn}
Both diagnostics run on the scalar log-likelihood trace only --- a cheap, convenient
proxy, \textbf{not} a substitute for a proper per-parameter convergence check. Treat a
value far from $1.0$ (or a small effective sample size relative to
\yamlkey{num\_iters}) as a reason to look closer with a dedicated tool
(\textsc{arviz}, \code{emcee}'s own diagnostics) before trusting the posterior shape.
\end{jwarn}

\section{Summary}
\label{sec:toymcmc-summary}

\sampler{ToyMCMC} is the family's reference method: independent Metropolis--Hastings
chains on the async pipeline, with no evidence estimate and no cross-chain coupling.
Table~\ref{tab:toymcmc-at-a-glance} gathers it.

\begin{table}[ht]
  \footnotesize
  \begingroup
  \setlength{\tabcolsep}{0pt}
  \begin{tabular}{@{}p{0.34\linewidth}p{0.66\linewidth}@{}}
    \toprule
    \textbf{Field} & \textbf{\sampler{ToyMCMC}} \\
    \midrule
    Kind & feedback (fixed budget: chains $\times$ iterations) \\
    Pipeline & async independent chains (\S\ref{sec:mcmc-runtime}) \\
    Point count & exactly \yamlkey{num\_chains}$\,\times\,$\yamlkey{num\_iters} \\
    Dimensions & any \\
    Required keys & \yamlkey{num\_chains}, \yamlkey{num\_iters},
      \yamlkey{proposal\_scale} \\
    u$\to$x mapping & yes \\
    \yamlkey{Bounds.seed} role & every chain's RNG stream and every proposal's
      \textsc{uuid} \\
    \textsc{uuid}s & deterministic from \code{(seed, chain\_id, step, stage)} \\
    Mid-run resume & yes --- time-cadenced, checkpoint validated against the card \\
    Extra output & \file{sampler\_summary.json}, \file{chain\_history.csv} \\
    \bottomrule
  \end{tabular}
  \endgroup
  \caption{\sampler{ToyMCMC} at a glance.}
  \label{tab:toymcmc-at-a-glance}
\end{table}

\begin{jtip}
Ready for a posterior that might trap a single chain in one mode? Chapter~\ref{ch:ptmcmc}
assumes everything in this chapter and covers only what parallel tempering adds.
\end{jtip}
